\section{Introducción}\label{sec:intro}

El problema de las $n$ reinas consiste en colocar $n$ reinas en un tablero de
$n\times n$ de modo que ninguna ataque a otra. En una solución hay exactamente una
reina en cada fila y en cada columna, y a lo sumo una en cada diagonal de cualquiera
de las dos direcciones. Denotamos por $Q(n)$ el número de soluciones.

Consideremos el tablero un paso antes de completarse. Si ya hay $n-1$ reinas que
no se atacan, quedan una sola fila y una sola columna vacías. Su intersección es
la única casilla en la que cabría la última reina sin mover las anteriores. La
colocación se completa si y sólo si ambas diagonales de esa casilla están libres.
Si elegimos uniformemente la colocación de las $n-1$ reinas, ¿con qué probabilidad
ocurre esto cuando crece el tablero?

Usaremos tableros pequeños como ejemplos. En las figuras numeramos filas y columnas
desde cero y escribimos primero la fila. La Figura~\ref{fig:ejemplos} muestra una
solución de $4\times4$ y otra de $5\times5$.

\begin{figure}[htbp]
\centering
\begin{tikzpicture}[tablero A]
 \begin{scope}
  \cuadriculaA{4}\trazarRejillaA{4}
  \reinaA{0}{1}\reinaA{1}{3}\reinaA{2}{0}\reinaA{3}{2}
  \panelA{2,5.34}{(a) Solución de $4\times4$}
 \end{scope}
 \begin{scope}[shift={(7,0)}]
  \cuadriculaA{5}\trazarRejillaA{5}
  \reinaA{0}{1}\reinaA{1}{3}\reinaA{2}{0}\reinaA{3}{2}
  \aroA{4}{4}\reinaA[colorAro]{4}{4}
  \panelA{2.5,5.34}{(b) Solución de $5\times5$}
 \end{scope}
\end{tikzpicture}
\caption{Dos soluciones pequeñas. Cada corona representa una reina; el aro señala la reina
$(4,4)$, que se usa en los ejemplos siguientes.}
\label{fig:ejemplos}
\end{figure}


En todo el artículo fijamos el tamaño del tablero y agrupamos las colocaciones según
el número de reinas que faltan: éste es el sentido de \emph{capa} en el título. Por
ejemplo, si en la Figura~\ref{fig:ejemplos}(b) retiramos la reina $(4,4)$, obtenemos
una colocación de la capa $h=1$ de $5\times5$; al colocarla de nuevo pasamos a la
capa $h=0$.

La misma pregunta tiene sentido cuando faltan $h$ reinas; diremos entonces que la
colocación tiene $h$ \emph{huecos}. Llamamos \emph{capa $h$}
al conjunto $\Om_h(n)$ de colocaciones legales de $n-h$ reinas. Sean
$N_h(n)=|\Om_h(n)|$ y $\mu_h$ la ley uniforme en esa capa cuando contiene al menos
una colocación. Tomamos una colocación con ley $\mu_h$. Sus $h$ filas vacías y sus
$h$ columnas vacías se cortan en $h^2$ \emph{cruces vacíos}; cuando hablemos de la posición de los
huecos nos referiremos a la de estos cruces. Elegimos uno de ellos
uniformemente. Diremos que
está \emph{disponible} si podemos añadir allí una reina, lo que equivale a que sus dos
diagonales estén libres. Para $h>1$, este experimento pregunta por la siguiente
inserción, aunque todavía queden reinas por colocar.

La probabilidad de una inserción legal se puede expresar mediante los tamaños de dos
capas consecutivas. Cada colocación de la capa $h-1$ tiene $n-h+1$ reinas que
podemos retirar para volver a la capa $h$. A la inversa, cada colocación de la
capa $h$ ofrece $h^2$ cruces para intentar la inserción. Al contar las inserciones
legales en ambos sentidos (Lema~\ref{lem:interpretacion}) se ve que la fracción media de
cruces disponibles bajo $\mu_h$ es el cociente
\begin{equation}\label{eq:fh}
 f_h(n)=\frac{(n-h+1)\,N_{h-1}(n)}{h^2\,N_h(n)} ,
\end{equation}
que tomamos como definición de $f_h(n)$ para $1\le h\le n$. Así, $f_1(n)$ responde
exactamente a la pregunta sobre la última reina. Demostramos que, para cada $h$ fijo,
$f_h(n)$ tiende a una misma constante positiva cuando crece el tablero.

Para ver qué se promedia, volvamos a la Figura~\ref{fig:ejemplos}(b) y
retiremos las reinas $(1,3)$ y $(4,4)$. Al quedar tres reinas en $5\times5$
tenemos dos huecos y cuatro cruces (Figura~\ref{fig:cruces-vacios}). Las filas
vacías son $1$ y $4$ y las columnas vacías son $3$ y $4$. La reina de $(3,2)$
ataca $(1,4)$ por la antidiagonal de suma $5$, y $(4,3)$ por la diagonal de
diferencia $1$. Los otros dos cruces admiten una reina.

\begin{figure}[htbp]
\centering
\begin{tikzpicture}[tablero A]
 \cuadriculaA{5}
 \foreach \rA in {1,4}{\fill[colorAro!7] (0,\rA) rectangle ++(5,1);}
 \foreach \cA in {3,4}{\fill[colorAro!7] (\cA,0) rectangle ++(1,5);}
 \foreach \rA in {1,4}{\foreach \cA in {3,4}{\fill[colorFondo] (\cA,\rA) rectangle ++(1,1);}}
 \trazarRejillaA{5}
 \ataqueA{3}{2}{1}{4}\ataqueA{3}{2}{4}{3}
 \reinaA{0}{1}\reinaA{2}{0}\reinaA{3}{2}
 \aroA{1}{3}\aroA{4}{4}\aspaA{1}{4}\aspaA{4}{3}
 \begin{scope}[shift={(6.0,.23)}]
  \node[anchor=west,font=\small\sffamily\bfseries] at (0,.25) {Los cuatro cruces};
  \aroA{1}{0}
  \node[anchor=west,font=\small] at (1.06,1.5) {Disponible: $(1,3)$ y $(4,4)$};
  \aspaA{2}{0}
  \node[anchor=west,font=\small] at (1.06,2.5) {Atacado: $(1,4)$ y $(4,3)$};
  \node[anchor=west,font=\small] at (.1,3.95) {$X_2(P)=2$ de $2^2=4$ cruces};
 \end{scope}
\end{tikzpicture}
\caption{Una colocación de la capa $\Omega_2(5)$. Los círculos y las aspas
ocupan las cuatro intersecciones de las filas y columnas vacías. Las líneas
discontinuas muestran por qué dos de ellas quedan bajo ataque.}
\label{fig:cruces-vacios}
\end{figure}


Escribimos $X_h(P)$ para el número de cruces disponibles en una colocación $P$.
En el ejemplo $X_2(P)=2$, de modo que un cruce elegido al azar en este tablero
admite una reina con probabilidad $2/4$. Para obtener $f_2(5)$ hay que promediar
esa fracción sobre todas las colocaciones legales de tres reinas en $5\times5$.

\subsection{Resultados principales}

Para identificar la constante, escribimos $x=h/n$ para la fracción de huecos. Al crecer $n$
con $h$ fijo, $x$ se acerca a cero. Esto lleva a estudiar el extremo de tablero lleno del perfil
variacional: la densidad $\nu_x$ en $[0,1]^2$, de masa $1-x$, que maximiza el funcional de
entropía de la Sección~\ref{subsec:variacional}. Para $x\in[0,1/2]$ tiene la forma
\[
 \nu_x(r,c)=b_x\,r_{R,x}(r)\,r_{C,x}(c)\,u_{D,x}(r-c)\,u_{A,x}(r+c),
 \qquad \int_0^1 r_{R,x}=\int_0^1 r_{C,x}=1,
\]
donde $r_{R,x}$ y $r_{C,x}$ describen los perfiles reescalados de filas y columnas
vacías, $u_{D,x}$ y $u_{A,x}$ las \emph{vacancias} de diagonales y antidiagonales (la densidad
de las que quedan vacías), y $b_x>0$ es la constante de normalización. Las propiedades del perfil se
demuestran en~\cite{Perfil}. Definimos
\begin{equation}\label{eq:F}
 F(x)=\frac{1-x}{b_x}=\int_{[0,1]^2}r_{R,x}(r)\,r_{C,x}(c)\,u_{D,x}(r-c)\,u_{A,x}(r+c)\,dr\,dc .
\end{equation}
Heurísticamente, $F(x)$ es la probabilidad de que un cruce vacío típico tenga libres sus dos
diagonales cuando hay $xn$ huecos: la posición del cruce se distribuye según $r_{R,x}r_{C,x}$ y
cada diagonal está vacía con densidad $u_{D,x}$ o $u_{A,x}$. Para las últimas capas la constante
relevante es $F(0)=1/b_0$.

\begin{teorema}[Ley de las últimas capas]\label{thm:limite}
Para cada $h\ge1$ fijo,
\[
 \lim_{n\to\infty}f_h(n)=F(0)=\frac1{b_0}\in(0,\infty).
\]
En particular, $\max_{1\le h\le K}|f_h(n)-F(0)|\to0$ para cada $K$ fijo.
\end{teorema}

Cada factor $f_j(n)$ relaciona dos capas consecutivas. Al multiplicar esas relaciones
desde $j=1$ hasta $j=h$, aparece el número de soluciones completas en el denominador.

\begin{corolario}[Cociente entre capas]\label{cor:cociente}
Para cada $h\ge1$ fijo,
\[
 \frac{N_h(n)}{Q(n)}=\frac{(n)_h}{(h!)^2\prod_{j=1}^hf_j(n)}
 \sim\frac{n^h}{(h!)^2F(0)^h},\qquad (n)_h=n(n-1)\cdots(n-h+1).
\]
\end{corolario}

El paso de la capa $h$ a la capa $h-1$ lleva la fracción de huecos de $h/n$ a
$(h-1)/n$. El corolario siguiente compara cada $f_h(n)$ con el valor de $F$
en el punto medio, $(h-\tfrac12)/n$.

\begin{corolario}[Las capas fijas no dejan constante residual]\label{cor:capas}
Para cada $K$ fijo,
\[
 \sum_{h=1}^K\log\frac{f_h(n)}{F\bigl((h-\tfrac12)/n\bigr)}\longrightarrow0,
\]
y lo mismo vale si $F((h-\tfrac12)/n)$ se sustituye por $F(h/n)$.
\end{corolario}

El número $K$ se fija antes de que crezca el tablero. Identificar el término
constante de $\log Q(n)$ exigiría controlar además las capas cuyo índice crece
con $n$. Para las dos últimas capas damos además una prueba directa de una cota inferior
explícita para $\liminf_nf_h(n)$, que no necesita la concentración de la cuenta inversa
(Sección~\ref{sec:inversa}).

\begin{teorema}[Cota inferior con constante explícita]\label{thm:positiva}
Sea $\kappa\in(0,1]$ la cota inferior de las vacancias diagonales $u_{D,x}$ y $u_{A,x}$ del perfil
(Teorema~\ref{thm:perfil}) y $a_0=(\kappa/2)^4/16$. Entonces
\[
 \liminf_{n\to\infty}f_h(n)\ge\frac{a_0}{2}>0,\qquad h=1,2.
\]
En particular, existen $c>0$ y $n_0$ tales que $f_1(n),f_2(n)\ge c$ para todo $n\ge n_0$.
\end{teorema}
